\section{The original infinitesimal monodromy and a flat logarithmic connection}
\label{sec:infinitesimal-braid}
This section differentiates the full integral braid of
Section~\ref{sec:full-braiding} without changing the original source
basis, full central weights, or either colour. Put $q=e^h$ in the
formal ring $\mathbb Q[[h]]$. The original Laurent matrices define
formal matrices on $X_{\mathbb Q}\otimes\mathbb Q[[h]]$, where
$X_{\mathbb Q}=X_A/(q-1)X_A\otimes\mathbb Q$ and its ordered basis
is the specialization of the original $1260$ vectors. Write
$E_c^0,F_c^0$ for the specialized matrices, and put
\[
 H_c b_i=(\lambda^c_{i1}-\lambda^c_{i2})b_i,\qquad
 D_c b_i=(\lambda^c_{i1}+\lambda^c_{i2})b_i=15b_i
 \quad(c=V,W).
\]
Thus $[H_c,E_c^0]=2E_c^0$, $[H_c,F_c^0]=-2F_c^0$ and
$[E_c^0,F_c^0]=H_c$. These follow by differentiating the original
Cartan relations and specializing the finite Laurent quantum integer
in the original commutator. Distinct colours commute. No original
matrix entry is replaced by its absolute value.

\subsection{Derivative in the fixed original coordinates}
The diagonal operator in the integral braid has expansion
\[
 Q(e^h)=1+hG+O(h^2),\qquad
 G=\sum_{c=V,W}\bigl(\Lambda_{c1}\otimes\Lambda_{c1}
                       +\Lambda_{c2}\otimes\Lambda_{c2}\bigr),
 \quad \Lambda_{ca}b_i=\lambda^c_{ia}b_i.
\]
Since $e^h-e^{-h}=2h+O(h^3)$, the $k=1$ term of each nilpotent
factor is $2hF_c^0\otimes E_c^0+O(h^2)$; every term with $k\ge2$
is divisible by $h^2$. The derivatives of the Chevalley matrices
in this term contribute only at order at least two. Consequently,
with the original flip $\tau$,
\begin{equation}
 R(e^h)=1+hr+O(h^2),\qquad c(e^h)=\tau(1+hr)+O(h^2),\qquad
 r=G+2\sum_c F_c^0\otimes E_c^0.
 \label{eq:ib-r}
\end{equation}
The monodromy here means the literal square of the same endomorphism
of $X_A\otimes X_A$. Multiplication in these fixed coordinates gives
\begin{align}
 c(e^h)^2&=1+h\Omega+O(h^2),\label{eq:ib-monodromy}\\
 \Omega&=r+\tau r\tau\nonumber\\
 &=\sum_{c=V,W}\bigl(D_c\otimes D_c+H_c\otimes H_c
       +2E_c^0\otimes F_c^0+2F_c^0\otimes E_c^0\bigr).
 \label{eq:ib-omega}
\end{align}
Indeed $2(\Lambda_1\otimes\Lambda_1+
\Lambda_2\otimes\Lambda_2)=D\otimes D+H\otimes H$.
Both central summands are $225$ times the identity, so their total
is $450I$. This scalar is part of the actual operator.

\subsection{Casimir identity and the precise sign}
For a single colour define the original classical action's Casimir
\[
 \mathcal C=H^2+2EF+2FE=H^2+2H+4FE.
\]
The primitive classical coproduct gives, by direct expansion,
\begin{equation}
 \Delta\mathcal C-\mathcal C\otimes1-1\otimes\mathcal C
 =2H\otimes H+4E\otimes F+4F\otimes E.
 \label{eq:ib-casimir}
\end{equation}
For completeness, it commutes with $H$, since $EF$ and $FE$ do.
Using the displayed form $H^2+2H+4FE$,
\[
 [\mathcal C,E]=2HE+2EH+4E-4HE
              =2(EH-HE)+4E=0,
\]
\[
 [\mathcal C,F]=-2HF-2FH-4F+4FH
              =2(FH-HF)-4F=0.
\]
On a highest vector $v$ of weight $m$, $Ev=0$ and hence
$\mathcal Cv=m(m+2)v$; commutation with $F$ gives that scalar
on every divided-power descendant.

Take actual highest source types $(m,p)$ and $(n,u)$, with
$m,p,n,u\in\{1,3,5,7,9\}$. Their first-colour tensor string of
highest weight $\ell=m+n-2r$ and second-colour string of highest
weight $t=p+u-2s$ have
$0\le r\le\min(m,n)$ and $0\le s\le\min(p,u)$.
The complete classical Clebsch maps are the $q=1$ specializations
of the maps proved in Section~\ref{sec:full-braiding}; the source
map is the unchanged actual $\Phi_D$ specialized at $q=1$.
Equations~\eqref{eq:ib-omega}--\eqref{eq:ib-casimir} therefore give
\begin{align}
 \omega_{m,n,r;p,u,s}
 &=450+\tfrac12\bigl(\ell(\ell+2)-m(m+2)-n(n+2)\bigr)
       \nonumber\\
 &\hspace{14mm}+\tfrac12\bigl(t(t+2)-p(p+2)-u(u+2)\bigr)
       \nonumber\\
 &=450+mn+pu-2r(m+n-r+1)-2s(p+u-s+1).
 \label{eq:ib-eigenvalue}
\end{align}
It acts identically on the retained multiplicity spaces. This is
also exactly the exponent of the full double braid, as follows
by multiplying the two explicitly proved highest-string scalars;
the two quantum-factorial ratios cancel and their powers of $q$
add to \eqref{eq:ib-eigenvalue}.

Every one-colour summand $mn-2r(m+n-r+1)$ is at least $-99$.
To prove this without dropping a range condition, interchange
$m,n$ if necessary so $m\ge n$. The difference of its values at
$r+1$ and $r$ is $-2(m+n-2r)<0$ for $0\le r<n$.
Its minimum is at $r=n$, where it is $-n(m+2)\ge-9\cdot11$.
Thus every eigenvalue of the actual $\Omega$ is at least $252$.
The complete actual-type spectrum below sharpens this uniform bound.

Here positivity can be stated for a specified positive form,
without using the original coordinate matrix as an unstated metric.
Choose any real basis in each actual highest multiplicity space
and declare it orthonormal. On a classical divided-power string
$e_i=F^{(i)}e_0$ of highest weight $m$, declare
\[
 \langle e_i,e_j\rangle=\delta_{ij}\binom mi.
\]
This is positive definite. The identities
$Fe_i=(i+1)e_{i+1}$ and $Ee_{i+1}=(m-i)e_i$ give
$(i+1)\binom m{i+1}=(m-i)\binom mi$, so $E^*=F$;
$H^*=H$ and $D^*=D$. Use the product form on both colours
and transport its orthogonal direct sum through the exact actual
map $\Phi_1$ to $X_{\mathbb R}$. This explicitly specifies the
metric on the original vector space. On its tensor square,
\eqref{eq:ib-omega} is self-adjoint. The proved complete string
decomposition and eigenvalue bound show
\[
 \langle v,\Omega v\rangle\ge252\langle v,v\rangle
 \qquad(v\in X_{\mathbb R}\otimes X_{\mathbb R}).
\]
A negative coefficient in an original source generator therefore
has not produced a negative infinitesimal-monodromy eigenvalue.
The exact transport giving this conclusion is
\eqref{eq:ib-r}--\eqref{eq:ib-eigenvalue} and the specified metric;
it does not assert that the generator coefficient is absent.

\subsection{A connection on the configuration space, with flatness proved}
For $N\ge2$ put
\[
 \operatorname{Conf}_N(\mathbb C)
 =\{(z_1,\ldots,z_N)\in\mathbb C^N:z_i\ne z_j\ (i\ne j)\},
 \qquad \alpha_{ij}=\frac{dz_i-dz_j}{z_i-z_j}.
\]
Let $\Omega_{ij}$ be the operator \eqref{eq:ib-omega} on positions
$i,j$ of the unchanged tensor product $X_{\mathbb C}^{\otimes N}$,
with identities elsewhere. On the trivial vector bundle with this
fibre define, for each specified constant $\kappa\in\mathbb C$,
\begin{equation}
 \nabla_\kappa=d-\kappa\sum_{1\le i<j\le N}\Omega_{ij}\alpha_{ij}.
 \label{eq:ib-connection}
\end{equation}
The following computation proves flatness on precisely this domain.
Disjoint pairs commute because their factors act in disjoint
positions. For either colour, direct commutation gives
\[
 [\Omega_{12},H_1+H_2]=[\Omega_{12},E_1+E_2]
 =[\Omega_{12},F_1+F_2]=[\Omega_{12},D_1+D_2]=0.
\]
For example, in the $E$ identity the Cartan term contributes
$2E_1H_2+2H_1E_2$, whereas the two root terms contribute
$-2E_1H_2-2H_1E_2$. The $F$ identity contributes
$-2F_1H_2-2H_1F_2$ and its negative. For $H$, each root
term has opposite weights in the two positions; central terms
commute. These prove every displayed identity, with both colours
adding independently. In particular
\begin{equation}
 [\Omega_{12},\Omega_{13}+\Omega_{23}]=0,
 \qquad [\Omega_{13},\Omega_{12}+\Omega_{23}]=0,
 \label{eq:ib-kohno}
\end{equation}
since the second operator in the first commutator is the sum,
over the two colours, of
\[
 (D_1+D_2)D_3+(H_1+H_2)H_3
       +2(E_1+E_2)F_3+2(F_1+F_2)E_3.
\]
The second identity is the first with indices permuted.
The one-forms are closed. For a triple set
$x=z_i-z_j$, $y=z_j-z_k$, so $z_i-z_k=x+y$. Direct wedge
multiplication gives
\begin{align*}
 \alpha_{ij}\wedge\alpha_{ik}
 -\alpha_{ij}\wedge\alpha_{jk}
 +\alpha_{ik}\wedge\alpha_{jk}
 &=\left(\frac1{x(x+y)}-\frac1{xy}
                  +\frac1{(x+y)y}\right)dx\wedge dy=0.
\end{align*}
If $T=[\Omega_{ij},\Omega_{ik}]$, the two commutator identities
in \eqref{eq:ib-kohno} make the three coefficients in the curvature
respectively $T,-T,T$. Their entire triple contribution is zero
by this wedge identity. Disjoint pairs contribute zero as proved,
and these exhaust all pairs of summands. Thus
$\nabla_\kappa^2=0$.

This construction relates the original full-source braiding to an
explicit flat differential connection by its formal derivative and
the proved infinitesimal braid relations. It does not identify the
analytic monodromy of \eqref{eq:ib-connection} with the original
quantum braid for finite $q$: such an identification is not used
in any proof here. No finite-field Frobenius realization or
canonical nonstandard plethysm geometry is supplied by this
characteristic-zero connection.
